\section{The Holy Night: First Vespers, Matins, and Midnight Mass}\label{sec:midnight}

At None on 24~December, the Vigil of the Nativity closes Advent. When the bells
ring for First Vespers that evening, Christmastide (\textit{Tempus Nativitatis},
General Rubric~72~a) begins. White vestments replace violet; the long waiting of
the prophets reaches the threshold of the night in which the Savior is born.

\subsection{First Vespers and the Night Watches of Matins}

First Vespers sets the pace of the evening with five antiphons that walk
step by step up to the stable door: \textit{Rex pacificus magnificatus est},
``King Peaceful exceeded all the kings of the earth''; \textit{Completi sunt
dies Mariae}, ``The days of Mary were accomplished that she should bring forth
her first-born Son''; and \textit{Levate capita vestra}, ``Lift up your heads:
behold, your redemption draweth nigh.'' After the short chapter from Titus
(3:4--5) rises the Christmas Vespers and Matins hymn \textit{Jesu Redemptor
omnium}, ``Jesus, the Ransomer of man,'' whose closing doxology---\textit{Jesu,
tibi sit gloria, qui natus es de Virgine}, ``To Jesus, from a Virgin sprung, be
glory given''---closes every iambic dimeter hymn of the Roman Breviary through
5~January. At the \textit{Magnificat}, the antiphon looks toward daybreak:
\textit{Cum ortus fuerit sol de coelo}, ``Or ever the sun be risen, ye shall see
the King of kings coming forth from the Father, as a bridegroom out of his
chamber.'' And at Compline, \textit{Alma Redemptoris Mater} drops its
Advent versicle for the Christmas cry: \textit{Post partum, Virgo, inviolata
permansisti / Dei Genitrix intercede pro nobis}, ``After thy delivery thou still
remainest a Virgin undefiled / Mother of God, pray for us,'' followed by the
Collect \textit{Deus, qui salutis aeternae} (Appendix~\ref{app:prayers}).

Earlier on the morning of 24~December, the choir at Prime had already set the
birth of Jesus inside the span of recorded history through the solemn
proclamation from the \textit{Roman Martyrology}: ``Upon the 25th day of
December... JESUS Christ, Himself Eternal God and Son of the Eternal Father,
being pleased to hallow the world by His most gracious coming, having been
conceived of the Holy Ghost, and when nine months were passed after His
conception, was born of the Virgin Mary at Bethlehem of Juda made Man.''

In the night hours before Midnight Mass, choral churches sing the three nocturns
of Christmas Matins. The Invitatory repeats the one sentence that governs the
night---\textit{Christus natus est nobis: Venite adoremus}, ``Unto us a Christ
is born: O come, let us worship Him.'' In the First Nocturn, the Church reads
three passages of Isaiah without announcement or title: Isaiah~9:1--6 (``For a
CHILD IS BORN to us, and a son is given to us''), Isaiah~40:1--8 (``Be
comforted, be comforted, my people, saith your God''), and Isaiah~52:1--6
(``Arise, arise, put on thy strength, O Sion''). The opening responsory
\textit{Hodie nobis coelorum Rex} marvels that today the King of heaven deigned
to be born of a Virgin to call lost man back to the heavenly kingdom, while the
second responsory sings \textit{Hodie per totum mundum melliflui facti sunt
coeli}, ``This day throughout the whole world the skies drop down sweetness.''

In the Second Nocturn, Pope St.~Leo the Great preaches his First Sermon on the
Nativity (Sermon~21), proclaiming that ``it would be unlawful to be sad to-day,
for to-day is Life's Birthday,'' and that Mary conceived Christ in her heart by
faith before she conceived him in her body. Between those lessons rise two of
the most beloved responsories in Western worship: \textit{O magnum mysterium},
``How great is this mystery, how wonderful is the teaching of the faith! The
beasts saw the new-born Lord lying in a manger,'' and \textit{Sancta et
immaculata virginitas}, ``O Mary, how holy and how spotless is thy
virginity---I am too dull to praise thee! For thou hast borne in thy breast Him
Whom the heavens cannot contain.'' In the Third Nocturn, the Breviary previews
all three Gospels of Christmas Day with three Latin Fathers: Pope St.~Gregory
the Great on Luke~2:1, St.~Ambrose on Luke~2:15, and St.~Augustine on John~1:1.
Then the \textit{Te Deum} rings out, and in traditional choral practice Midnight
Mass begins at once.

\subsection{Why Three Masses on One Day}

In the \textit{Summa Theologiae} (Tertia Pars, q.~83, a.~2, ad~2), St.~Thomas
Aquinas explains why the Church appoints three distinct Masses on 25~December:

\begin{quote}
On Christmas Day, however, several masses are said on account of Christ's
threefold nativity. Of these the first is His eternal birth, which is hidden in
our regard, and therefore one mass is sung in the night, in the ``Introit'' of
which we say: ``The Lord said unto Me: Thou art My Son, this day have I begotten
Thee.'' The second is His nativity in time, and the spiritual birth, whereby
Christ rises ``as the day-star in our hearts'' (2~Pet.\ 1:19), and on this
account the mass is sung at dawn, and in the ``Introit'' we say: ``The light
will shine on us today.'' The third is Christ's temporal and bodily birth,
according as He went forth from the virginal womb, becoming visible to us
through being clothed with flesh: and on that account the third mass is sung in
broad daylight, in the ``Introit'' of which we say: ``A child is born to us.''
\end{quote}

Aquinas immediately notes the complementary reading as well: in the flesh,
Christ was literally born at night, entering the darkness of our weakness, so
the Midnight Mass reads the Gospel of his temporal birth (Luke~2:1--14);
meanwhile his eternal generation belongs to unapproachable light, so the Day
Mass reads St.~John's Prologue (John~1:1--14). Together the three Masses knit
eternity, history, and the believer's soul into one celebration.

\subsection{The First Mass at Midnight}

The Roman station is \textit{ad Sanctam Mariam maiorem ad Praesepe}, St.~Mary
Major at the Crib. As the priest approaches the altar, the choir sings the
Introit from Psalm~2:7: \textit{Dominus dixit ad me: Filius meus es tu, ego
hodie genui te}, ``The Lord hath said to me: Thou art my son, this day have I
begotten thee.'' After the \textit{Kyrie}, the \textit{Gloria in excelsis Deo}
returns for the first time since the end of November, taking up on earth the
angels' own hymn from Bethlehem.

\begin{massbox}{Collect of the First Mass at Midnight}
\textbf{Latin} (1962 Roman Missal). Deus, qui hanc sacratissimam noctem veri
luminis fecisti illustratione clarescere: da, quaesumus; ut, cujus lucis
mysteria in terra cognovimus, ejus quoque gaudiis in caelo perfruamur: Qui
tecum vivit et regnat in unitate.

\textbf{English} (Lasance, 1945). O God, Who hast brightened this most holy
night with the shining of the true light, grant, we beseech Thee, that we may
enjoy in heaven the delights of Him whose mystical light we have known on earth.
Who with Thee.
\end{massbox}

The Collect moves from faith to beatific vision: tonight we know the mysteries
of the true Light on earth by faith, asking that we may enjoy his delights in
heaven face to face. The Epistle (Tit~2:11--15) immediately tells the assembly
how those who know that light must live: ``For the grace of God our Saviour hath
appeared to all men: Instructing us, that, denying ungodliness and worldly
desires, we should live soberly and justly and godly in this world, Looking for
the blessed hope and coming of the glory of the great God and our Saviour Jesus
Christ.'' Preaching on these verses in his fifth Homily on Titus, St.~John
Chrysostom notes that St.~Paul joins Christ's two appearings---the first in
grace, the second in glory---and deliberately chooses a strong verb: ``He has
not said `avoiding,' but `denying.' Denying implies the greatest distance, the
greatest hatred and aversion.''

The Gradual and Communion both draw from Psalm~109~[110]:3---\textit{Tecum
principium in die virtutis tuae: in splendoribus sanctorum, ex utero ante
luciferum genui te}, ``With thee is the principality in the day of thy strength:
in the brightness of the saints: from the womb before the day star I begot
thee''---anchoring the infant's cry in the Father's timeless begetting before
the morning star.

Then the deacon or priest chants the Gospel of Luke~2:1--14: Augustus's census
of the whole world, Mary bringing forth her firstborn son, wrapping him in
swaddling clothes and laying him in a manger because there was no room in the
inn, and the angel breaking the silence over the shepherds watching their
flocks: ``Fear not; for, behold, I bring you good tidings of great joy that
shall be to all the people: For, this day is born to you a Saviour, who is
Christ the Lord, in the city of David'' (Luke~2:10--11). In Homily~8 (read at
Matins just before this Mass), Pope St.~Gregory the Great points straight to
those first lines: Caesar enrolled the world in time to show that One was coming
in the flesh who enrolls his elect in eternity, and he was born in
Bethlehem---``the House of Bread''---because he is the Living Bread come down
from heaven to feed the faithful.

In his \textit{Discourse for Christmas Night} (\textit{The Incarnation, Birth
and Infancy of Jesus Christ}), St.~Alphonsus Maria de Liguori speaks directly to
every sinner kneeling in the dark nave:

\begin{quote}
And what tidings could be a greater joy to a nation of poor exiles, condemned to
death, than for them to be told that their Saviour had come, not only to deliver
them from death, but also to obtain for them permission to return to their
country? And this is what I announce to you this night: \textit{A Saviour is
born to you!} Jesus Christ is born; and he is born for you, to deliver you from
everlasting death, and to open to you heaven, which is our country, and from
which we had been banished in punishment for our sins.
\end{quote}

At the Creed, all kneel at the words \textit{Et incarnatus est de Spiritu Sancto
ex Maria Virgine: et homo factus est}---``And was incarnate by the Holy Ghost of
the Virgin Mary, and was made man.'' In the Secret (\textit{Accepta tibi sit}), the Church names the altar's offering
\textit{per haec sacrosancta commercia}, ``through this sacred intercourse,''
asking that we be found made like unto him in whom our human substance is
already united with God.

Then the Preface of the Nativity opens the Canon with a sentence that governs
the entire season through 5~January---\textit{Quia per incarnati Verbi
mysterium, nova mentis nostrae oculis lux tuae claritatis infulsit: ut dum
visibiliter Deum cognoscimus, per hunc in invisibilium amorem rapiamur}---declaring
that through the mystery of the incarnate Word a new light of God's brightness
has shone upon the eyes of our mind, so that knowing God visibly we may be
caught up through him into the love of things invisible. And at this Mass alone
in the entire year, the proper \textit{Communicantes} inside the Canon speaks of
the night itself: \textit{Communicantes, et noctem sacratissimam celebrantes},
``keeping that most holy night, in which the spotless virginity of the blessed
Mary brought forth a Saviour to this world.'' At Dawn and Day, and through
1~January, that word will change to \textit{diem sacratissimum}, ``keeping that
most holy day.'' Only here, in the quiet of midnight, does the Roman Canon hold
the night itself in its hands.
